\chapter{\xxx{Optimization methods}}\label{chap:optimization-methods}

As discussed in~\cref{chap:code-optimization}, code optimization is crucial for achieving efficient execution of programs on modern hardware architectures such as multi-core CPUs and GPUs.
In this chapter, we outline several prominent optimization methods and frameworks that have been developed to facilitate (automated) code optimization in high-performance computing applications, putting particular emphasis on those that are relevant to our work --- targeting data layout and traversal optimizations.

Code optimization methods can target different points in the compilation and execution pipeline.
% Some methods operate above the source code level, providing high-level abstractions or domain-specific languages (DSLs) that allow developers to express computations in a way that is amenable to more powerful optimizations. Other methods perform optimizations at the source code level, possibly using annotations or pragmas to guide the optimization process.
Methods that operate above the source code level often provide high-level abstractions or domain-specific languages (DSLs) that allow developers to express computations in a way that is amenable to more powerful optimizations. Methods that perform optimizations at the source code level may use annotations or pragma directives to guide the optimization process without changing the syntax or semantics of the programming language, or use libraries that provide the necessary functionality for optimization while still allowing developers to write code in a familiar way.

More low-level methods operate on intermediate representations (IRs) or directly on the generated machine code, applying transformations to improve performance based on static analysis or empirical measurements.

Some methods may also involve runtime systems that adaptively optimize code during execution based on observed performance characteristics; in such cases, methods like just-in-time (JIT) compilation techniques may be employed to generate optimized code on-the-fly based on current execution patterns.

We primarily focus on abstractions and frameworks that operate above or at the source code level, as these are most relevant to our work. However, we also briefly discuss the wider landscape of optimization methods to provide context for our contributions.


\section{Compiler optimizations}\label{sec:compiler-optimizations}

Compiler optimizations are a fundamental aspect of modern compilers. In high-performance computing, compilers such as LLVM~\cite{lattner2004llvm} and GCC~\cite{gcc} implement a wide range of optimization techniques to improve the performance of source code, starting from basic optimizations that remove abstractions and redundant computations, inline functions, perform sub-expression elimination, etc. More advanced optimizations include loop transformations (e.g., loop unrolling, loop fusion, invariant code motion), vectorization, and inter-procedural optimizations that analyze and optimize across function boundaries.

In the context of typical high-performance (scientific computing) applications with quite simple control flow and heavy numerical computations, these optimizations mainly allow using higher-level abstractions (e.g., \cpp templates) without significant performance penalties (most notably, when using so-called zero-cost abstractions~\cite{stroustrup2012foundations}). They can also optimize arithmetic expressions, eliminate redundant computations, and perform some basic loop transformations. However, these optimizations by themselves are often insufficient to achieve optimal performance on modern hardware architectures, especially for complex loop nests and data access patterns.

% Some compilers and frameworks, such as MLIR~\cite{lattner2021mlir}, provide extensible intermediate representations that allow custom representation and optimization passes to be implemented --- some targeting specific domains or hardware architectures. This is particularly useful for domain-specific languages (DSLs) that can leverage MLIR to implement their own optimization pipelines, including polyhedral optimizations, autotuning, and other transformations not typically part of general-purpose compilers. While MLIR is built on top of traditional static-single-assignment (SSA)-based IRs\footnote{In an SSA-based IR, each variable is assigned exactly once and thus a sequence of computations unbroken by control flow can be represented as a directed acyclic graph (DAG) of data dependencies.}, it allows defining custom dialects and facilitates boilerplate such as parsing, code generation, and so on.

Since compilers are designed to be general-purpose, they often rely on heuristics and cost models that may not capture the specific performance characteristics of specific applications or hardware architectures. This can be partially mitigated by using special compiler flags or using profile-guided optimizations (PGO) that allow the compiler to make more informed decisions based on empirical performance data. However, standard compiler optimizations still often fail to achieve optimal performance even with these methods.


\subsection{Polyhedral optimization}\label{sec:polyhedral-optimization}

Workloads comprising nested loops with affine bounds and affine memory accesses (i.e., Static Control Parts, or SCoPs) are common in high-performance computing applications (see, e.g., \citeauthor{bacon1994compiler}~\cite{bacon1994compiler}). To target such applications, \citeauthor{feautrier1992some}~\cite{feautrier1992some} and others pioneered the use of the polyhedral model to represent and optimize computations in loop nests. Polyhedral optimization frameworks, such as \textbf{Polly}~\cite{grosser2012polly} (in LLVM), \textbf{GRAPHITE}~\cite{trifunovic2010graphite} (in GCC), and \textbf{PLuTo}~\cite{bondhugula2008pluto} (source-to-source), use the powerful mathematical framework of the polyhedral model (provided by libraries such as ISL~\cite{verdoolaege2010isl}) to represent computations in loops in a more compact and analyzable form compared to traditional compiler IRs such as Static Single Assignment (SSA) forms.
% Polyhedra offer a way to represent data layouts, loop iteration spaces, and dependencies between computations in a unified manner.

Polyhedral optimization is the application of integer linear algebra to optimize loop nests with affine bounds and affine memory accesses. The main idea is representing the iteration space of loops, as well as memory accesses and dependencies, using polyhedra to enable, for example, finding hyperplanes orthogonal to data dependencies that can represent parallelizable loops, or finding transformations of the iteration space that improve data locality or enable vectorization.

There are works that apply the theory of polyhedra and integer linear programming to optimize data layouts as well. For example, the work by \citeauthor{clauss2000automatic}~\cite{clauss2000automatic} uses Ehrhart polynomials to change accesses to arrays into consecutive memory accesses, optimizing the data layout for the target cache hierarchy.

Many works iterate on the basic polyhedral model to extend its applicability or improve the optimization results by combining it with cost models, heuristics, or machine learning techniques (e.g., \citeauthor{baghdadi2021deep}~\cite{baghdadi2021deep}) to evaluate transformations proposed by the polyhedral model and find good parameters for these transformations (e.g., tiling sizes). Some works show benefits of combining polyhedral optimization based on cost models or heuristics with autotuning~\cite{chelini2020automatic} or human insights~\cite{zinenko2014clint,bagneres2016opening}.

Some DSL compilers (e.g., Halide~\cite{ragan2013halide}, Tiramisu~\cite{baghdadi2019tiramisu}) that will be discussed in~\cref{sec:abstractions} use polyhedral optimization techniques (or derived methods) as part of their optimization pipelines. These DSLs often provide higher-level abstractions that allow developers to guide the optimization process using directives that can be mapped to polyhedral transformations.

\xxx{\textbf{Polygeist}~\cite{moses2021polygeist} is a tool that translates \cpp code to MLIR (Multi-Level Intermediate Representation), which is a modern compiler infrastructure that provides a flexible and extensible framework for representing and optimizing computations. Polygeist uses the Clang frontend to parse C/\cpp source code and then raises the resulting low-level LLVM IR to higher-level MLIR representations, particularly its affine dialect. This allows Polygeist to apply MLIR's polyhedral optimization passes (e.g., affine fusion, loop tiling, and parallelization) to originally imperative C/\cpp code. It combines these polyhedral optimizations with traditional SSA-based compiler optimizations, outperforming both Polly and PLuTo on the PolyBench~\cite{pouchet2012polybench} benchmark suite.}

Since polyhedral optimizations are powerful tools that can fundamentally restructure the code, some works have focused on improving the interpretability of the transformations and their results, such as the work by \citeauthor{bagneres2016opening}~\cite{bagneres2016opening} that translates the scheduling decisions of polyhedral optimizers into a human-readable set of transformation directives --- e.g., loop interchange and strip-mine (tiling) --- that describe the transformations in a more intuitive way.


\section{Autotuning}\label{sec:autotuning}

Autotuning techniques are powerful methods for optimizing performance of applications based on empirical performance measurements or cost models. They involve performance-cost model-based and empirical search methods for exploring large parameter spaces (\emph{search space}s) to find optimal configurations. They can control various parameters, such as algorithmic choices, data layouts, tiling sizes, unrolling factors, and parallelization strategies, or even tune build configurations and compiler flags.

Many search algorithms have been proposed for exploring the search space, performing local or global optimization. Simple methods such as grid search and random search can be effective for search spaces with a small number of parameters. Hill climbing methods greedily explore the search space by iteratively moving towards better configurations based on performance measurements. These methods are surprisingly effective in many cases, especially when the search space is not too large or when the performance landscape does not have many local minima. Even methods that simply combine random search with local search (e.g., by performing random restarts) can be quite effective in practice.

More sophisticated methods include simulated annealing, genetic algorithms, swarm optimization, Bayesian optimization, and many others. Many of these methods balance exploration and exploitation in different ways, and their effectiveness can depend on the specific characteristics of the search space and the performance landscape.

\textbf{OpenTuner}~\cite{ansel2014opentuner} is the first framework that combines multiple search techniques (e.g., genetic algorithms, simulated annealing, etc.) in ensembles (based on multi-armed bandits) with a flexible API for defining the search space and the objective function. It is a general-purpose autotuning framework in Python that allows users to define their own search spaces and objective functions and perform black-box evaluations of the target application with different configurations from the search space.

% \textbf{ATF}~\cite{rasch2018atf} is a \cpp framework that builds on top of OpenTuner and a domain-specific autotuner CLTune. It provides options for pruning the search space by defining constraints on the parameters (e.g., divisibility constraints for tiling sizes) as well as employing parallelism when generating the search space.

% Tools such as \textbf{Hyperopt}~\cite{bergstra2015hyperopt} and \textbf{Optuna}~\cite{akiba2019optuna} are general-purpose optimization frameworks (originally built for optimizing hyperparameters of machine learning models) that use Bayesian optimization techniques to efficiently explore large search spaces with a limited number of evaluations. Bayesian optimization involves building a surrogate model of the performance landscape based on observed measurements, and using this model to guide the search towards promising regions of the search space.

% These tools can be used for autotuning by defining the search space and objective function in a way suitable for optimizing code performance. Optuna stands out for its define-by-run API, which allows for dynamic construction of the search space during the optimization process. This can be particularly useful for complex optimization problems where the search space may not be known a priori, cannot be efficiently represented in a static manner, or is too large to be exhaustively explored.

% There are many autotuning frameworks that are specific to certain domains, such as AutoTVM~\cite{chen2018tvm} for deep learning workloads, Kernel Tuner~\cite{van2019kernel} for GPU kernels

Autotuning techniques can be used in combination with other optimization methods, such as polyhedral optimization or DSLs\@. They are particularly useful for optimizing performance in cases where the performance landscape is complex and difficult to model analytically.


\section{Abstractions}\label{sec:abstractions}

Flexible abstractions allowing manipulation of data layouts and iteration orders have been the subject of extensive development \xxx{(cite)}. They provide a means to optimize performance for different hardware architectures without needing to change the core algorithmic logic. These abstractions can also serve as flexible API layers for libraries (such as linear algebra libraries~\cite{whaley1998automatically}).

We can broadly categorize these abstractions into the following groups that we will discuss in more detail in the following sections:

\begin{description}
  \item[Domain-specific languages (DSLs):]\label{des:dsl} DSLs provide high-level abstractions for expressing computations in a way that is amenable to powerful optimizations (often using custom optimization pipelines built for the specific domain).
  \item[Annotation-based tools:]\label{des:annotation-based-tools} Annotation-based tools allow developers to provide hints or directives in the source code (e.g., using pragmas or attributes) that guide the optimization process. Perhaps the most well-known example of this is OpenMP~\cite{dagum1998openmp}, which is a standard widely supported by most compilers that allows developers to annotate loops and sections of code for parallel execution.
  \item[Parallelization frameworks and libraries:]\label{des:parallelization-frameworks} Parallelization frameworks and libraries provide APIs for expressing parallel computations in the source code without necessarily relying on custom compilation pipelines.
\end{description}


\subsection{Domain-Specific languages}\label{sec:dsls}

DSLs are a powerful tool for optimizing performance in specific domains. Some DSLs such as Halide~\cite{ragan2013halide} and Tiramisu~\cite{baghdadi2019tiramisu} have front-ends that are directly embedded in \cpp, allowing developers to write code in a familiar language while benefiting from the high-level abstractions. Other DSLs, such as Exo~\cite{ikarashi2022exocompilation} do not focus this aspect and instead provide ways to easily extend the framework with custom optimization transformations or specific hardware instructions. These three DSLs are particularly relevant to our work as they are designed for separation of algorithms from optimization decisions and they are used in works explore automatic optimizations or computer-aided algorithm design, which are the main themes of our contributions. A common feature of these DSLs is that they verify the correctness of transformations using the polyhedral (or derived) model, which is especially useful in optimization pipelines that speculatively apply random transformations (e.g., in autotuning).

\textbf{Halide} was primarily designed for image processing applications, but its abstractions and optimization techniques have been applied to more general stencil and tensor computations as well. Its design philosophy is based on the separation of algorithm and user-defined scheduling, allowing developers to express the algorithm in a high-level way while providing a set of scheduling directives that can be used to optimize the performance for different hardware architectures. However, its expressiveness is somewhat limited when it comes to more complex computations, such as those involving irregular data structures or imperfectly nested loops, making it not suitable for general computations involving cycles; as such, its compilation pipeline and the scheduling directives use a restricted version of the polyhedral model that specifically targets computations expressible as directed acyclic graphs of array computations. The design principles of Halide have been applied to works such as TVM~\cite{chen2018tvm} for deep learning workloads.

\textbf{Tiramisu} extends the expressiveness of Halide by allowing more complex computations. It targets image processing, deep learning, as well as general scientific computing applications. Unlike Halide, Tiramisu can express computations with irregular data structures and imperfectly nested loops, which allows it to be applied to much wider range of applications. Its compilation pipeline is based on the polyhedral model, and it provides a rich set of scheduling directives that allow developers to optimize performance for different hardware architectures.

\textbf{Exo} is a general-purpose DSL for exocompilation, which is the process of generating optimized code for specific hardware architectures. Unlike Halide and Tiramisu, Exo is defined as a Python library that provides a flexible API for defining computations and optimization transformations. Similarly to the other DSLs, it allows developers to express computations in a high-level way while providing a set of optimization transformations that can be used to optimize performance for different hardware architectures. However, where it stands out is in its extensibility, allowing developers to easily add custom optimization transformations or specific hardware instructions that are defined intuitively as Python functions specifying their semantics.


\subsection{Annotation-based approaches}\label{sec:annotation-based}

Works grouped in this category are perhaps the least aligned with the approaches explored in our contributions. However, they too provide powerful tools for separating algorithmic logic from optimization decisions. Their main advantage is that they operate directly on the source code without changing the syntax or semantics of the language.

The most well-known example of this is \textbf{OpenMP}~\cite{dagum1998openmp}, which is a standard widely supported by most compilers that allows developers (or automatic tools) to annotate loops and sections of code with pragma directives that indicate how the code should be parallelized. OpenMP provides a rich set of directives for expressing different types of parallelism (e.g., data parallelism, task parallelism) and for controlling various aspects of the parallel execution (e.g., scheduling, synchronization, privatization).

Some efforts have been made to extend OpenMP with directives for more complex, composable transformations. One example is the work by \citeauthor{kruse2018user}~\cite{kruse2018user} that extends the OpenMP API in the Clang compiler with directives for loop transformations such as tiling, reversal (the loop is executed in reverse order), interchange, and array packing (some subset of the accessed array is copied to a temporary buffer).

Works such as \textbf{Loopy}~\cite{namjoshi2016loopy} provide annotation-based directives for expressing algorithm transformations in a simple math-like syntax. For example, if a loop has iteration variables $i$, $j$, and $k$ representing nested loop indices, the directive \texttt{affine(LoopLabel, {[i, j, k] -> [i, k, j]})} would indicate interchange of the nested loops. It stands out for its verification capabilities, allowing developers to verify the correctness of transformations via the polyhedral model.


\subsection{Parallelization frameworks and libraries}\label{sec:parallelization-frameworks}

The category of parallelization frameworks and libraries is quite broad, and can be further subdivided into frameworks that provide high-level APIs for expressing parallel computations (e.g., Kokkos~\cite{trott2021kokkos}, RAJA~\cite{beckingsale2019raja}, MPI~\cite{mpi-5.0}, GridTools~\cite{afanasyev2021gridtools}) and libraries that provide native abstractions usable in the source code (e.g., mdspan~\cite{trott2022mdspan}, CuTe~\cite{thakkar2023cutlass}). Note that some of the frameworks in the first group provide (usually native) abstractions for data layouts that are of particular relevance to our work, such as \lstinline{Kokkos::View} and \lstinline{RAJA::View} for multidimensional arrays.
% We chose examples of frameworks and libraries that have a non-trivial overlap with the approaches explored in our contributions and that are widely used in the high-performance computing community.

\textbf{Mdspan} is a \cpp standard abstraction for multidimensional views of data. It provides a flexible and efficient way to model the mapping of multidimensional array indices to linear memory addresses with arbitrary strides. Specifically, the mapping is defined by a layout policy that typically represents an affine function that maps the input index tuple to a linear offset in memory (non-affine mappings are possible, but require manual handling). Mdspan also allows for expressing subviews (e.g., slices, subarrays) and supports accessing data elements using custom load and store semantics. However, the abstraction does not provide any built-in mechanisms for expressing transformations of data layouts or defining layouts as compositions of simpler layouts.

\textbf{Kokkos} and \textbf{RAJA} are \cpp libraries that provide high-level APIs for expressing parallel computations in a way that is portable across different hardware architectures (e.g., multi-core CPUs, GPUs). They provide abstractions for parallel loops, reductions, and other common parallel patterns that are useful in a wide range of HPC applications. They also provide native abstractions for multidimensional arrays (Kokkos::View and RAJA::View) similar to mdspan, which also facilitate memory management and data transfers between different memory spaces (e.g., host and device memory). However, their abstractions are not designed to allow flexible manipulation of data layouts and iteration orders, which is a key aspect of our contributions.

\textbf{MPI} is a widely used library for distributed computing that provides a rich set of APIs for communication, synchronization, collective operations, remote memory access, and other features that are essential for building scalable parallel applications. MPI provides an opaque abstraction for data layouts (MPI datatypes) that is flexible enough to represent a wide range of data layouts, including non-contiguous and otherwise complex layouts. The flexibility of MPI datatypes, combined with MPI being a widely supported standard, makes it an interesting framework in the context of our contributions, which we explore in more detail in \cref{sec:mpi-bindings}.

\textbf{GridTools} is a \cpp library designed to facilitate the development of stencil computations on structured grids. Its abstractions are native to the \cpp/CUDA source code, and it enables developers to express the computations in a way that allows for powerful optimizations, including data caching and parallelization.

\textbf{CuTe} is a set of \cpp abstractions in the Cutlass library, primarily designed for expressing tensor computations on NVIDIA GPUs\@. It provides a flexible way to express data layouts as well as defining an algebra of layouts that serves as a powerful tool for defining various derived layouts (e.g., tiled layouts). Similarly to mdspan, a CuTe layout is defined as an affine mapping from the multidimensional index space to linear memory addresses, but CuTe allows each index dimension to be replaced with a tuple of new dimensions whose combined size is the same as the original dimension (this can represent tiling, for example); side-effect of this property is that each CuTe layout has a well-defined linearization (can be indexed by a single linear index). The algebra of layouts is defined by functional composition of layouts (output of one layout used to index the other). CuTe is perhaps the closest to our goal of defining layout-agnostic algorithms as its layouts separate the index space from the physical memory layout (a feature shared with mdspans and Kokkos/RAJA views), and it also supports some transformations of layouts (e.g., tiling) via its layout algebra.


\section{Machine learning and LLM-directed optimization}\label{sec:ml-directed-optimization}

Recently, machine learning (ML) techniques, including large language models (LLMs), have been increasingly applied to various aspects of code optimization. The main motivation for using these techniques is their flexibility and ability to learn and generalize complex patterns from data, which can overcome limitations of traditional optimization methods that rely on hand-crafted heuristics or analytical models.

% We have already discussed autotuning techniques in~\cref{sec:autotuning} that can be used to tune various compiler flags and options based on empirical measurements, so here we focus on ML-based methods that are not better categorized as autotuning. What characterizes these methods is that they make decisions based on learned patterns, applying them to new code or new optimization problems without relying on empirical measurements of the target code, which would be the case for autotuning techniques (such as Bayesian optimization methods).

\subsection{ML-based surrogate models and directive generation}\label{sec:ml-surrogate-models}

Some older works have applied ML techniques as surrogate models for performance prediction, such as the work by \citeauthor{fursin2011milepost}~\cite{fursin2011milepost} in compiler optimization. Recent additions in the area of HPC include works, such as \citeauthor{baghdadi2021deep}~\cite{baghdadi2021deep}, who use deep learning methods as cost models in a polyhedral optimization framework (the Tiramisu compiler), or \citeauthor{adams2019learning}~\cite{adams2019learning}, who use a learned cost model to guide the search in the Halide autotuning framework.

More recently, LLMs have emerged as a powerful tool for code optimization, with the ability to learn from large codebases and generate optimized code based on natural language descriptions or code snippets. Various works have explored the use of off-the-shelf LLMs or fine-tuned (or even trained from scratch) LLMs for generating directives used to optimize code, such as OpenMP directives for parallelization, scheduling directives for frameworks such as Halide or Tiramisu, or even generating optimized code directly.

Works that use LLMs for generating optimization directives include PragFormer~\cite{harel2025pragformer} that trains a transformer-based model to identify \lstinline{for} loops in C code that can be parallelized and even identify the need for privatization or reduction variables. Works such as OMPify~\cite{kadosh2023advising}, Graph2Par~\cite{chen2023learning}, and AutoParLLM~\cite{mahmud2023autoparllm} then use GNN-based (Graph Neural Network) models to generate OpenMP directives for parallelization with even higher accuracy.

HPC-Coder~\cite{nichols2024hpc} shows that fine-tuning LLMs on HPC code can enable them to generate OpenMP directives for parallelization with high accuracy as well as generate parallel code or MPI-based distributed code directly. The MPI code generation is further studied by works such as MPI-RICAL~\cite{schneider2023mpi} and MPI{rigen}~\cite{schneider2024mpirigen}.

Works that use LLMs or other deep learning methods for generating scheduling directives include ComPilot~\cite{merouani2025agentic}, which uses an LLM agent as a search driver when exploring the search space of transformation directives, and LOOPer~\cite{merouani2025looper}, which uses an LSTM-based solution inspired by the work of \citeauthor{baghdadi2021deep}. Both solutions achieve autoscheduling in the Tiramisu compiler.

% Using ML techniques as surrogate models or requesting them to propose optimization directives has the benefit of enabling formal verification of the optimized code, as the model acts simply as a search driver in a solution that would otherwise work as an autotuning solution, trying semi-random configurations and adjusting heuristics for selecting each next configuration based on observed performance measurements.

\subsection{Direct generation of optimized code}\label{sec:ml-direct-code-generation}

Using ML techniques as surrogate models or requesting them to propose optimization directives has the benefit of enabling formal verification, as the model acts simply as a search driver. However, direct generation of optimized code goes further and allows the model to make more fundamental changes to the code. This was not a viable option until recently, as the proficiency of LLMs in generating optimized code was not sufficient for this task.

% Some solutions even use the machine learning model to generate the entire optimized code directly from a natural language description (or from another high-level representation of the expected behavior of the code)~\cite{godoy2024large}.

Since LLMs are generally trained on large codebases, and only a handful of generally available code falls into the HPC domain, the performance of off-the-shelf LLMs on generating optimized HPC code has started to become competitive only recently, perhaps with the introduction of Codex~\cite{chen2021evaluating}. This challenge was earlier addressed by fine-tuning LLMs on HPC code~\cite{nichols2024performance} or by using smaller models trained from scratch on HPC code~\cite{xu2022systematic}. Such work has to be accompanied by standardized benchmarks and datasets of HPC code, such as PolyBench~\cite{pouchet2012polybench}, ParEval~\cite{nichols2024can}, KernelBench~\cite{ouyang2025kernelbench}, HPCorpus~\cite{kadosh2023quantifying}, or HeCBench~\cite{jin2023benchmark}.

Since the proficiency of LLMs as of the beginning of 2026 is still not sufficient for generating optimized code without supervision or verification, most works rely on using LLMs for translating between a more interpretable language and a programming model such as CUDA~\cite{palkowski2024gpt,dearing2024lassi}. One of these works, LASSI, defines a self-correcting pipeline where the LLM proposes a solution similarly to an autotuning solution, with the added benefit of being able to respond to feedback from the compiler or the verifier. This direction is used in various state-of-the-art solutions such as ComPilot~\cite{merouani2025agentic}.
